\documentclass[12pt,a4paper]{article}

\usepackage{amssymb}
\usepackage{amsmath}
\usepackage{amsthm}
\usepackage{bm}
\usepackage{enumitem}
\usepackage{mathtools}
\usepackage{lipsum}
\usepackage[brazil]{babel}
\usepackage[utf8]{inputenc}
\usepackage[T1]{fontenc}
\usepackage[pdftex]{graphicx}
\usepackage{color}
\usepackage{tikz}
\usetikzlibrary{arrows, calc, external}
\usepackage{verbatim}
\usepackage{url}
\usepackage{longtable}
\usepackage{geometry}
\geometry{
    a4paper,
    margin=2cm
}

\usepackage{listings}
\usepackage{xcolor}

\definecolor{codegreen}{rgb}{0,0.6,0}
\definecolor{codegray}{rgb}{0.5,0.5,0.5}
\definecolor{codepurple}{rgb}{0.58,0,0.82}
\definecolor{backcolour}{rgb}{0.9,0.9,0.9}

\lstset{
    % Formatação
    language=C++,
    basicstyle=\ttfamily\small,
    tabsize=4,
    keepspaces=true,
    columns=fullflexible,
    lineskip=-2pt,
    % Cores
    backgroundcolor=\color{backcolour},
    commentstyle=\color{codegreen},
    keywordstyle=\color{blue}\bfseries,
    stringstyle=\color{codepurple},
    % Numeração
    numbers=left,
    numberstyle=\tiny\color{codegray},
    numbersep=6pt,
    % Quebras
    breaklines=true,
    breakatwhitespace=false,
    % Invisíveis
    showspaces=false,
    showstringspaces=false,
    showtabs=false,
    % Layout
    frame=single,
    rulecolor=\color{black},
    captionpos=b
}

\usepackage{helvet}
\renewcommand{\familydefault}{\sfdefault}

\usepackage{hyperref}
%\usepackage[alf]{abntex2cite}
%\providecommand{\abntnextkey}{}

\usepackage{titlesec}
\usepackage{setspace}

\newtheorem{prop}{Proposição}

\begin{document}
\onehalfspacing

\begin{center}
    \textbf{Fundação Universidade Federal do ABC}\\
    \textbf{Pró-reitoria de pesquisa}\\
    \small Av. dos Estados, 5001, Santa Terezinha, Santo André/SP, CEP 09210-580\\
    Bloco L, 3º Andar, Fone (11) 3356-7617 | iniciacao@ufabc.edu.br
\end{center}

\vspace{0.8cm}

\noindent \textbf{Relatório Parcial de Iniciação Científica referente ao Edital:} 01/2025 PIC\slash PIBIC\slash PIBITI\slash PIBIC-AF \\[0.5cm]
\noindent \textbf{Nome do aluno:} Gabriel Frigo \\[0.5cm]
\noindent \textbf{Assinatura do aluno:} \hrulefill \\[0.5cm]
\noindent \textbf{Nome do orientador:} Cristiane Maria Sato \\[0.5cm]
\noindent \textbf{Assinatura do orientador:} \hrulefill \\[0.5cm]
\noindent \textbf{Título do projeto:} Problemas de Fluxos em Redes: Teoria, Algoritmos e Implementações \\[0.5cm]
\noindent \textbf{Palavras-chave do projeto:} Otimização Combinatória, Teoria dos Grafos, Fluxos \\[0.5cm]
\noindent \textbf{Área do conhecimento do projeto:} Ciência da Computação, Matemática Discreta \\[0.5cm]
\noindent \textbf{Bolsista:} ( \textbf{\textsf{X}} ) Sim \quad ( \quad ) Não. \\
\textbf{Em caso positivo, informar a modalidade da bolsa:} PIBIC - CNPq \\[1cm]

\begin{center}
    \textbf{Santo André}\\
    \textbf{\today}
\end{center}

\newpage
\renewcommand{\contentsname}{Sumário}
\tableofcontents
\newpage

\section{Resumo}
Fluxos em digrafos constituem uma área fundamental da Otimização Combinatória, com amplas aplicações práticas em redes de transporte e alocação de recursos. Este projeto adota uma abordagem que combina fundamentos teóricos, implementação algorítmica e validação experimental para o estudo de problemas de fluxo máximo e fluxo de custo mínimo (veja~\cite{ahuja1993network}). Nesta etapa parcial, foram estudados algoritmos de fluxo máximo e técnicas de modelagem. A validação prática está sendo realizada através de experimentos computacionais utilizando instâncias de \textit{benchmarks} reconhecidas (veja~\cite{Jensen2022}).

\section{Introdução}
A Otimização Combinatória é uma área de pesquisa situada na interseção entre Matemática e Ciência da Computação (veja~\cite{korte2018combinatorial}). Do ponto de vista matemático, a área mantém forte influência com a Teoria dos Grafos, uma vez que digrafos constituem objetos combinatórios com amplo poder de modelagem (veja~\cite{bondy2008graph, diestel2017graph}).
Sob a perspectiva computacional, o foco reside na obtenção de algoritmos eficientes que produzam soluções garantidamente ótimas (veja~\cite{schrijver2003combinatorial, cook1998combinatorial}). Dentre os tópicos de destaque, os problemas de fluxos em digrafos ocupam posição central, tratando de redes onde recursos devem ser transportados respeitando restrições de capacidade (veja~\cite{ahuja1993network}).
O objetivo deste projeto é proporcionar uma formação sólida na área, integrando teoria e prática. Até o presente momento, as atividades focaram na consolidação teórica e na implementação de algoritmos para o problema do fluxo máximo.

\section{Fundamentação teórica}
O problema do $st$-fluxo máximo, onde se busca maximizar a quantidade de fluxo enviada de um vértice fonte $s$ para um vértice de destino $t$, é um dos pilares deste estudo (veja~\cite{ford1956maximal}). Foram analisados algoritmos como os de Ford-Fulkerson, Edmonds-Karp e Dinic (veja~\cite{ford1956maximal, edmonds1972theoretical, dinic1970algorithm}).
O projeto também contempla o problema de fluxo de custo mínimo, que busca minimizar o custo total enquanto satisfaz as demandas da rede. O algoritmo principal a ser implementado para este fim é o \textit{network simplex}, uma especialização do algoritmo simplex (veja~\cite{dantzig1951application, ahuja1993network}).
A fundamentação teórica revisada nesta etapa reafirma que muitos problemas combinatórios, como o emparelhamento máximo em digrafos bipartidos, podem ser resolvidos via reduções para problemas de fluxo (veja~\cite{ahuja1993network}).

\section{Metodologia}
\subsection{Materiais e Métodos}
A metodologia baseia-se em ciclos de estudo bibliográfico seguidos de implementação. Os algoritmos estão sendo desenvolvidos em linguagem C++ para garantir alta performance, utilizando listas de adjacência otimizadas.
Os experimentos serão conduzidos em instâncias padronizadas das coleções DIMACS (veja~\cite{Jensen2022}), permitindo a comparação de desempenho. A análise baseia-se em métricas de tempo de execução e corretude dos algoritmos implementados frente às soluções ótimas conhecidas (veja~\cite{ford1956maximal}).

\subsection{Etapas da pesquisa}
As atividades do projeto foram estruturadas e encontram-se no seguinte status:
\begin{itemize}
    \item \textbf{Atividade 1:} Estudo dos fundamentos de Otimização Combinatória e Teoria dos Grafos (veja~\cite{cook1998combinatorial}). \textbf{[Concluída]}
    \item \textbf{Atividade 2:} Estudo teórico do problema de fluxo máximo e algoritmos clássicos (veja~\cite{ford1956maximal, edmonds1972theoretical, dinic1970algorithm, goldberg1988new}). \textbf{[Concluída]}
    \item \textbf{Atividade 3:} Implementação em C++ dos algoritmos de fluxo máximo. \textbf{[Concluída]}
    \item \textbf{Atividade 4:} Estudo de técnicas de modelagem e redução de problemas. \textbf{[Concluída]}
    \item \textbf{Atividade 5:} Implementação de problemas modelados (Emparelhamento e Caminhos Disjuntos). \textbf{[Concluída]}
    \item \textbf{Atividade 6:} Elaboração do relatório parcial (este documento). \textbf{[Em andamento]}
    \item \textbf{Atividade 7:} Estudo teórico do problema de fluxo de custo mínimo (veja~\cite{cook1998combinatorial}). \textbf{[A iniciar]}
    \item \textbf{Atividade 8:} Implementação em C++ do algoritmo \textit{network simplex}. \textbf{[A iniciar]}
    \item \textbf{Atividade 9:} Testes experimentais em instâncias de larga escala. \textbf{[A iniciar]}
    \item \textbf{Atividade 10:} Análise comparativa de desempenho dos algoritmos. \textbf{[A iniciar]}
    \item \textbf{Atividade 11:} Elaboração do relatório final e artigo de conclusão. \textbf{[A iniciar]}
\end{itemize}

\section{Resultados parciais}

Nesta seção, apresentamos os avanços teóricos e práticos alcançados na primeira metade do projeto. A exposição é estruturada de forma a construir a fundamentação matemática elementar das redes de fluxo, evoluir para as propriedades de redes residuais, culminar nas implementações de algoritmos de estado-da-arte (\textit{Edmonds-Karp}, \textit{Dinic} e a família \textit{Push-Relabel}) e, por fim, demonstrar o poder destas abstrações na Otimização Combinatória por meio de reduções clássicas validadas empiricamente.

\subsection{Base Matemática dos Grafos de Fluxo}

Para a formulação rigorosa do problema de fluxo, baseamo-nos nas definições clássicas da Teoria dos Grafos e Otimização Combinatória (veja~\cite{ahuja1993network, bondy2008graph}). Uma rede de fluxo é formalmente definida como um digrafo (grafo direcionado) $D = (V, A)$, onde $V$ representa o conjunto de vértices e $A$ o conjunto de arcos.

A topologia desta rede é dotada de dois vértices terminais distintos: uma origem $s \in V$ (fonte ou \textit{source}) e um destino $t \in V$ (sumidouro ou \textit{sink}). Cada arco $a \in A$, que conecta um vértice $u$ a um vizinho $v$, é representado pelo par ordenado $(u,v)$ e possui associada uma capacidade não-negativa estipulada por uma função $c: A \rightarrow \mathbb{R}_{\ge 0}$. Esta capacidade $c(u,v)$ dita o limite superior do recurso (físico ou informacional) que pode transitar diretamente através desta conexão.

Um fluxo em $D$ é caracterizado por uma função $f: A \rightarrow \mathbb{R}_{\ge 0}$ que atribui a cada arco $(u,v)$ um valor de tráfego $f(u,v)$. Para que este roteamento seja classificado como um fluxo factível (ou viável) do ponto de vista matemático, a função $f$ deve satisfazer estritamente dois axiomas fundamentais (veja~\cite{ford1956maximal, cook1998combinatorial}):

\begin{enumerate}
    \item \textbf{Restrição de Capacidade:} O fluxo roteado em qualquer arco não pode exceder a sua capacidade estrutural intrínseca, tampouco assumir valores negativos. Formalmente:
    \begin{equation}
        0 \le f(u,v) \le c(u,v), \quad \forall (u,v) \in A
    \end{equation}

    \item \textbf{Lei de Conservação de Fluxo:} Com exceção dos vértices terminais ($s$ e $t$), nenhum vértice intermediário na rede possui a capacidade de gerar, consumir ou reter fluxo. A soma total do fluxo que ingressa em um vértice $v$ deve ser rigorosamente igual à soma do fluxo que dele emerge. Matematicamente, para todo $v \in V \setminus \{s, t\}$:
    \begin{equation}
        \sum_{\{u \mid (u,v) \in A\}} f(u,v) = \sum_{\{w \mid (v,w) \in A\}} f(v,w)
    \end{equation}
\end{enumerate}

O valor global do fluxo, comumente denotado por $|f|$, quantifica a rede em sua totalidade, representando a quantidade líquida de recurso que efetivamente deixa a fonte $s$ e, por consequência direta da lei de conservação, ingressa no sumidouro $t$. O objetivo central do clássico Problema do Fluxo Máximo consiste justamente em determinar uma função factível $f$ que maximize o escalar $|f|$.

Esta abstração matemática de rede, capacidades e restrições serve como o substrato exato para o encapsulamento orientado a objetos desenvolvido nas etapas práticas deste projeto. A classe-base C++ desenvolvida abstrai este formalismo garantindo, em nível de \textit{software}, que a invariante da restrição de capacidade nunca seja violada.

\subsection{Redes Residuais, Caminhos Aumentantes e Dualidade}

A busca sistemática pela maximização do fluxo em uma rede requer uma estrutura auxiliar que quantifique não apenas a capacidade ociosa das conexões, mas também a possibilidade de reversão de decisões de roteamento subótimas estabelecidas em iterações prévias. Para este fim, introduz-se o concept de digrafo residual (ou rede residual) $D_f = (V, A_f)$ (veja~\cite{ahuja1993network, ford1956maximal}).

Dado um fluxo viável $f$ operando sobre a rede original $D$, a capacidade residual para qualquer par ordenado de vértices $(u,v)$ é formalmente definida pela equação:
$$c_f(u,v) = c(u,v) - f(u,v) + f(v,u)$$.

Esta formulação matemática é de suma importância estrutural, pois o termo algébrico $+ f(v,u)$ abstrai a simetria reversa do fluxo, garantindo topologicamente a viabilidade de cancelar e redirecionar um fluxo previamente enviado no sentido oposto. O conjunto de arcos residuais $A_f$ é então constituído exclusivamente pelos pares $(u,v)$ cujas capacidades residuais satisfazem rigorosamente a condição $c_f(u,v) > 0$.

Operando sobre a topologia de $D_f$, define-se um caminho aumentante como qualquer caminho simples com origem na fonte $s$ e término no sumidouro $t$. A existência de tal caminho implica de forma direta que o fluxo global $|f|$ pode ser estritamente incrementado sem violar as restrições estruturais da rede original. O incremento máximo viável ao longo deste trajeto é delimitado pelo arco de menor capacidade residual, comumente denominado arco gargalo.

O critério de parada para os algoritmos de aumento iterativo, bem como o certificado matemático de otimalidade, é assegurado pelo Teorema do Fluxo Máximo-Corte Mínimo (\textit{Max-Flow Min-Cut Theorem}) (veja~\cite{ford1956maximal}). O teorema introduz a noção de um corte $s-t$, definido como uma partição do conjunto de vértices $V$ em dois subconjuntos mutuamente exclusivos $S$ e $T$, tais que $s \in S$ e $t \in T$. A capacidade deste corte é computada pela soma escalar das capacidades de todos os arcos originais direcionados de $S$ para~$T$.

O teorema demonstra uma dualidade perfeita em otimização: o valor do fluxo máximo em uma rede é estritamente igual à capacidade do corte mínimo. Computacionalmente, a ausência de um caminho aumentante na rede residual $D_f$ é a condição matemática necessária e suficiente para atestar que o fluxo atual atingiu o seu limite ótimo global (veja~\cite{cook1998combinatorial}).

\subsubsection{Prova do Teorema do Fluxo Máximo-Corte Mínimo}

A demonstração deste teorema é construtiva e apoia-se diretamente na topologia do digrafo residual (veja~\cite{ford1956maximal, cook1998combinatorial}). Seja $f$ um fluxo em $D$. O valor do fluxo $|f|$ é sempre menor ou igual à capacidade de qualquer corte $s-t$ na rede, uma vez que o tráfego de recursos é fisicamente limitado pelas conexões que particionam a fonte do sumidouro. Portanto, se encontrarmos um fluxo $|f|$ cujo valor seja exata e algebricamente igual à capacidade de algum corte específico, teremos provado simultaneamente que $f$ é máximo e que o corte em questão é o mínimo.

Suponha que $f$ seja um fluxo máximo. Por definição, não existe um caminho aumentante de $s$ para $t$ na rede residual $D_f$ (veja~\cite{ahuja1993network}). Definimos o conjunto $S$ como o conjunto de todos os vértices que são alcançáveis a partir da fonte $s$ em $D_f$ através de arcos com capacidade residual estritamente positiva ($c_f > 0$). Definimos $T = V \setminus S$ como o restante dos vértices.

Como não há caminho aumentante, o sumidouro $t$ é inalcançável a partir de $s$ em $D_f$, logo $t \in T$ e $s \in S$. Temos, portanto, um corte $s-t$ válido particionando a rede.

Para qualquer arco original $(u,v) \in A$ tal que $u \in S$ e $v \in T$, sabemos que $v$ não é alcançável a partir de $u$ em $D_f$. Isto implica obrigatoriamente que a capacidade residual $c_f(u,v) = 0$. Pela definição de capacidade residual, isto só ocorre se o arco original estiver completamente saturado, ou seja, $f(u,v) = c(u,v)$. De forma análoga, para qualquer arco no sentido reverso, de $T$ para $S$, o fluxo deve ser nulo ($f(v,u) = 0$), caso contrário haveria capacidade residual reversa permitindo a travessia de $S$ para $T$ (veja~\cite{cook1998combinatorial}).

Conclui-se que o fluxo líquido através do corte $(S, T)$ é exatamente a soma das capacidades dos arcos de $S$ para $T$, o que corresponde à capacidade total do corte. Como o fluxo da rede deve passar através deste gargalo, provamos que o valor $|f|$ iguala a capacidade deste corte, atestando a otimalidade de ambos (veja~\cite{ford1956maximal}).

\subsubsection{O Teorema da Integralidade}

Para que os problemas clássicos de Otimização Combinatória possam ser reduzidos a instâncias de fluxo, um último pilar matemático faz-se necessário: a garantia de que as soluções encontradas pertencerão ao domínio dos números inteiros ($\mathbb{Z}$).

O Teorema da Integralidade de Ford-Fulkerson estabelece que: se a função de capacidade $c(u,v)$ assume valores inteiros para todos os arcos $(u,v) \in A$, então existe um fluxo máximo $f$ onde o valor $f(u,v)$ transitando em cada arco é estritamente inteiro (veja~\cite{ford1956maximal}).

A prova assenta-se na natureza iterativa do método dos caminhos aumentantes. Iniciando com um fluxo nulo (inteiro), e sabendo que as capacidades são inteiras, o gargalo residual em qualquer iteração será sempre um valor inteiro (veja~\cite{ahuja1993network}). Consequentemente, as operações de adição e subtração durante a atualização do fluxo preservarão a integralidade da rede indefinidamente. Esta propriedade é a chave teórica que habilita o uso dos algoritmos desenvolvidos para a resolução de problemas de emparelhamento e de alocação discreta.

\subsection{Arquitetura Orientada a Objetos e Estruturas de Dados}

A transição dos modelos matemáticos abstratos para implementações computacionais de alto desempenho exige decisões arquiteturais rigorosas. Para garantir a manutenibilidade, a extensibilidade e a correção empírica dos experimentos, a biblioteca desenvolvida neste projeto foi fundamentada no paradigma da Orientação a Objetos através da linguagem C++ moderno.

Antes mesmo da definição estrutural, a implementação estabelece apelidos de tipo (\textit{Type Aliases}), como \lstinline{Long} e \lstinline{Size}, para garantir a semântica e a portabilidade do código entre diferentes arquiteturas. Outro cuidado fundamental em nível de \textit{hardware} foi a definição das constantes de limite (\lstinline{INF} e \lstinline{MAX}). Ao invés de utilizar o valor absoluto máximo suportado pelos tipos numéricos, aplicou-se um deslocamento de bits para a direita (\lstinline{>> 8}). Esta técnica atua como uma margem de segurança matemática, prevenindo que somas subsequentes de fluxos e capacidades resultem em comportamentos indefinidos de estouro de memória (\textit{Integer Overflow}).

O núcleo do sistema é regido pela classe base abstrata \lstinline{FlowNetwork}. Esta classe encapsula a representação topológica do digrafo e o estado iterativo do fluxo, isolando a complexidade estrutural das lógicas específicas de busca de cada algoritmo. A topologia é gerenciada pelas seguintes estruturas fundamentais:

\begin{itemize}
    \item \textbf{\lstinline{Edge}:} Uma estrutura (\textit{struct}) leve que representa um arco direcionado, agregando as variáveis \lstinline{from} (origem), \lstinline{to} (destino), \lstinline{capacity} (capacidade original $c$) e \lstinline{flow} (fluxo atual $f$).
    \item \textbf{\lstinline{size}:} Variável escalar que armazena a cardinalidade do conjunto de vértices ($|V|$), utilizada para dimensionar as estruturas auxiliares dinamicamente.
    \item \textbf{\lstinline{edges}:} Um vetor linear que armazena todos os arcos da rede. A alocação contígua em memória favorece o \textit{cache} do processador, mitigando os custos de latência comuns em alocações dinâmicas de digrafos.
    \item \textbf{\lstinline{adj}:} Uma lista de adjacências otimizada (\lstinline{std::vector<std::vector<Size>>}). Em vez de armazenar ponteiros ou cópias dos arcos, este vetor bidimensional guarda estritamente os índices numéricos que os arcos incidentes ocupam no vetor principal \lstinline{edges}, garantindo acesso em tempo constante $\mathcal{O}(1)$.
\end{itemize}

Para preservar a integridade matemática da rede residual e evitar a duplicação de código (\textit{DRY Principle - Don't Repeat Yourself}), a manipulação do fluxo foi estritamente confinada a dois métodos protegidos (\textit{protected}):

\begin{itemize}
    \item \textbf{\lstinline{get_residual_capacity(Size edge_id)}:} Função que computa em tempo real a capacidade residual $c_f(u,v)$ subtraindo o fluxo alocado da capacidade original. O uso do atributo \lstinline{[[nodiscard]]} garante em tempo de compilação que o valor computado não seja acidentalmente ignorado pelas rotinas de busca.
    \item \textbf{\lstinline{push_flow(Size edge_id, Long flow_amount)}:} Este é o método mais crítico da arquitetura, responsável por rotear fisicamente uma quantidade de fluxo pela rede. Ele explora um truque algorítmico basilar: ao invés de buscar o arco reverso em uma lista, o sistema garante no momento de inserção (\lstinline{add_edge}) que todo arco direto e seu respectivo arco reverso sejam alocados sequencialmente em pares (índices 0 e 1, 2 e 3, etc.). Consequentemente, o método utiliza a operação bit a bit Ou Exclusivo (\lstinline{id ^ 1ULL}) para acessar instantaneamente o arco complementar, somando o fluxo no arco direto e o subtraindo no arco reverso em uma única instrução de máquina.
\end{itemize}

No escopo da interface pública, para contornar as restrições de instanciação polimórfica inerentes ao C++, a arquitetura incorpora os padrões de projeto \textit{Factory Method} e \textit{Prototype} (veja~\cite{gamma1994design}) através dos métodos virtuais \lstinline{make} e \lstinline{clone}. O método \lstinline{make} atua como um construtor virtual dinâmico, permitindo a alocação de uma rede inteiramente nova e vazia que preserva a mesma classe algorítmica do objeto original (e.g., instanciar um novo \textit{Dinic} a partir de um ponteiro genérico), eliminando o acoplamento aos tipos concretos. Em complemento, o método \lstinline{clone}, de particular relevância para as aplicações deste projeto, garante a capacidade de gerar cópias profundas (\textit{deep copies}) da rede carregando todo o seu estado topológico e residual em tempo de execução. Esta funcionalidade estrutural do \lstinline{clone} é crucial para cenários algorítmicos complexos que exigem a simulação de múltiplos cortes ou roteamentos sobre a mesma topologia base, preservando o estado original intacto entre as iterações.

Por fim, a classe expõe a interface polimórfica primária \lstinline{compute_max_flow(Size source, Size sink)}, um método virtual puro que atua como um contrato estrutural obrigatório para todas as subclasses. Além disso, para viabilizar a extração de resultados práticos após a convergência do fluxo (como a identificação das conexões ativas em um Emparelhamento Bipartido Máximo), a classe provê o método \lstinline{get_edges()}. O retorno da estrutura através de uma referência constante (\lstinline{const &}) permite a leitura externa de toda a malha de arcos sem comprometer o rigor e a segurança do encapsulamento orientado a objetos. A consolidação estrutural de todas estas decisões arquiteturais e a definição completa da interface C++ encontram-se detalhadas no Código~\ref{lst:flownetwork} (disponível no Apêndice~\ref{ap:flownetwork}).

\subsection{Algoritmos Baseados em Caminhos Aumentantes}

A resolução computacional do problema de fluxo máximo fundamenta-se historicamente no método iterativo proposto por Ford e Fulkerson (veja~\cite{ford1956maximal}). A premissa estrutural deste paradigma consiste na manutenção de uma rede residual $A_f$ e na busca sistemática por caminhos aumentantes — percursos simples direcionados da fonte $s$ ao sumidouro $t$ onde todos os arcos possuem capacidade residual estritamente positiva.

O método estipula que, enquanto houver um caminho aumentante em $A_f$, o fluxo global da rede deve ser incrementado pelo valor da capacidade gargalo deste caminho (a menor capacidade residual encontrada no trajeto). Contudo, a definição matemática original estabelece apenas o critério de viabilidade topológica do percurso, não impondo nenhuma diretriz algorítmica ou heurística de seleção de qual caminho processar em digrafos onde múltiplos trajetos aumentantes coexistem (veja~\cite{ahuja1993network}).

Esta abstração geométrica, embora elegante para a prova do Teorema do Fluxo Máximo-Corte Mínimo, apresenta uma vulnerabilidade computacional severa. Em instâncias teóricas onde as capacidades dos arcos assumem valores irracionais, demonstra-se matematicamente que uma escolha arbitrária ou patológica de caminhos pode levar o método a realizar infinitas iterações. Em casos extremos, a sequência de aumentos não apenas falha em terminar, como o valor do fluxo pode convergir assintoticamente para um limite estritamente inferior ao fluxo máximo real da rede (veja~\cite{cook1998combinatorial}).

Mesmo no domínio prático das redes com capacidades inteiras, uma má escolha iterativa pode forçar o método a operar em tempo $\mathcal{O}(|A| \cdot |f^*|)$, onde $|f^*|$ é o valor escalar do fluxo máximo. Esta dependência direta do valor numérico restringe a abordagem a um tempo de execução pseudopolinomial, cujo desempenho degrada exponencialmente em função da magnitude das capacidades originais da rede.

Para transpor esta barreira analítica e garantir limites de complexidade de tempo estritamente polinomiais — independentes do valor numérico das capacidades ou do fluxo máximo —, a literatura em Otimização Combinatória evoluiu aplicando restrições rigorosas sobre as políticas de exploração da rede residual. As duas abordagens de caminhos aumentantes mais bem-sucedidas e clássicas implementadas nesta biblioteca, o Algoritmo de Edmonds-Karp e o Algoritmo de Dinic, nascem fundamentalmente da necessidade de disciplinar a seleção iterativa destes percursos.

\subsubsection{O Algoritmo de Edmonds-Karp}

A superação da complexidade pseudopolinomial do método genérico de Ford-Fulkerson foi alcançada de forma independente por Edmonds e Karp (veja~\cite{edmonds1972theoretical}) e por Dinic (veja~\cite{dinic1970algorithm}). O algoritmo de Edmonds-Karp concretiza a restrição topológica mais intuitiva e elegante: a cada iteração, o percurso selecionado para o roteamento deve ser rigorosamente o caminho aumentante mais curto da fonte $s$ ao sumidouro $t$ na rede residual $A_f$, assumindo-se que todos os arcos viáveis possuam comprimento unitário.

A implementação desta diretriz fundamenta-se na delegação da busca a uma Busca em Largura (BFS - \textit{Breadth-First Search}). O impacto teórico desta restrição é profundo, pois assegura a propriedade da monotonicidade das distâncias: ao longo da execução do algoritmo, a distância topológica mínima (mensurada em número de arcos) da fonte $s$ a qualquer vértice $v \in V$ na rede residual nunca decresce (veja~\cite{cook1998combinatorial}).

Para demonstrar a complexidade polinomial $\mathcal{O}(|V| \cdot |A|^2)$, analisa-se a frequência com que um arco pode figurar como o gargalo limitante do fluxo. Quando um arco $(u,v)$ atua como gargalo, ele é completamente saturado e removido de $A_f$. Por pertencer ao caminho mais curto, a relação de distâncias no momento da sua saturação é estritamente $d(v) = d(u) + 1$.

Para que o arco $(u,v)$ retorne à rede residual em uma iteração futura, o algoritmo deverá obrigatoriamente rotear fluxo na direção reversa $(v,u)$. Isto implica topologicamente que o trecho $(v,u)$ agora pertence a um novo caminho mais curto, exigindo que a nova distância da fonte até $u$ satisfaça a equação $d'(u) = d'(v) + 1$. Pela lei da monotonicidade, sabendo que a distância do vértice $v$ não pode ter diminuído ($d'(v) \ge d(v)$), deduz-se algebricamente que $d'(u) \ge d(u) + 2$ (veja~\cite{ahuja1993network}).

Conclui-se analiticamente que, entre duas saturações consecutivas de uma mesma conexão, a distância do vértice de origem em relação à fonte sofre um incremento de, no mínimo, duas unidades. Uma vez que o limite superior para a distância de um caminho simples em um digrafo é $|V| - 1$, cada arco da rede pode ser saturado no máximo $\mathcal{O}(|V|)$ vezes ao longo de toda a execução.

Existindo $\mathcal{O}(|A|)$ arcos na malha geométrica, o limite superior de aumentos globais é estritamente circunscrito a $\mathcal{O}(|V| \cdot |A|)$. Multiplicando este fator de repetição pelo custo computacional intrínseco de $\mathcal{O}(|A|)$ operações de uma busca BFS, o tempo de execução de pior caso é cravado rigorosamente em $\mathcal{O}(|V| \cdot |A|^2)$ (veja~\cite{korte2018combinatorial}). A materialização desta heurística, utilizando a herança e o encapsulamento orientado a objetos da interface principal, encontra-se estruturada no Código~\ref{lst:edmonds_karp} (disponível no Apêndice~\ref{ap:edmonds_karp}).

\subsubsection{O Algoritmo de Dinic}

Embora o algoritmo de Edmonds-Karp apresente um comportamento fortemente polinomial, o limite assintótico de $\mathcal{O}(|V| \cdot |E|^2)$ pode se mostrar proibitivo em instâncias massivas ou em digrafos de alta densidade estrutural. O algoritmo formulado por Yefim Dinitz (veja~\cite{dinic1970algorithm}) contorna esta barreira analítica ao introduzir o processamento em fases e a estruturação de uma abstração topológica denominada Digrafo de Níveis (\textit{Level Graph}).

A evolução central desta arquitetura baseia-se na separação rigorosa entre a descoberta de caminhos e o roteamento de fluxo. No início de cada fase iterativa, uma Busca em Largura (BFS) global é executada a partir da fonte $s$ para atribuir níveis topológicos $d(v)$ a todos os vértices, correspondentes às suas distâncias mínimas exatas na rede residual $A_f$. O Digrafo de Níveis gerado é, por definição, um subdigrafo direcionado acíclico (DAG) composto estritamente pelos arcos residuais que avançam exatamente um nível em direção ao sumidouro, satisfazendo a restrição geométrica $d(v) = d(u) + 1$ (veja~\cite{ahuja1993network}).

Uma vez validado e construído este subdigrafo acíclico, o algoritmo abandona a estratégia de incremento singular e emprega uma Busca em Profundidade (DFS - \textit{Depth-First Search}) para determinar o chamado Fluxo Bloqueador (\textit{Blocking Flow}). A DFS explora a topologia de forma ávida, saturando simultânea e exaustivamente múltiplos caminhos admissíveis presentes no Digrafo de Níveis, até que o sumidouro $t$ se torne inalcançável sob as restrições de nível vigentes (veja~\cite{cook1998combinatorial}).

Para assegurar a viabilidade temporal desta busca exaustiva, a implementação integra uma estrutura imperativa de otimização de estado denominada ponteiros mortos (\textit{Dead-end Pointers}). Mantendo um vetor na memória que registra a exata posição do último arco explorado a partir de cada vértice, o algoritmo inviabiliza que caminhos sem saída (\textit{dead-ends}) ou conexões com capacidade esgotada sejam revisitados durante a mesma fase.

Esta poda dramática na árvore de recursão garante uma amortização matemática rigorosa: a extração completa do fluxo bloqueador opera no limite estrito de $\mathcal{O}(|V| \cdot |A|)$ operações (veja~\cite{korte2018combinatorial}). Como a propriedade de monotonicidade garante que a distância de $s$ a $t$ aumente pelo menos uma unidade a cada novo Digrafo de Níveis construído, o sistema processa um limite superior absoluto de $|V|$ fases. Multiplicando o custo da DFS por fase pelo número máximo de fases, a complexidade assintótica global é inquestionavelmente estabelecida em $\mathcal{O}(|V|^2 \cdot |A|)$ (veja~\cite{dinic1970algorithm}). A materialização destas otimizações de estado na arquitetura C++ encontra-se no Código~\ref{lst:dinic} (disponível no Apêndice~\ref{ap:dinic}).

\subsection{Algoritmos de Prefluxo: Introdução}

Todos os algoritmos clássicos baseados no método de Ford-Fulkerson (incluindo Edmonds-Karp e Dinic) compartilham uma invariante fundamental: a lei de conservação de fluxo é rigorosamente mantida em todos os vértices intermediários ao longo de toda a execução. Sob essa ótica, os algoritmos operam mantendo um fluxo sempre viável e esforçam-se iterativamente para alcançar a otimalidade (a inexistência de caminhos aumentantes).

Em 1988, Andrew Goldberg e Robert Tarjan (veja~\cite{goldberg1988new}) propuseram uma inversão dramática deste paradigma. A família de algoritmos por eles introduzida inverte a invariante: a otimalidade (inexistência de caminhos da fonte ao sumidouro na rede residual) é mantida desde a configuração inicial, mas a viabilidade do fluxo é deliberadamente relaxada. A lei de conservação é suspensa durante as iterações intermediárias, permitindo que os vértices acumulem recursos temporariamente. Esta estrutura de roteamento relaxada é matematicamente definida como um pré-fluxo (\textit{preflow}) (veja~\cite{ahuja1993network}).

Formalmente, um pré-fluxo $f$ é uma função que respeita estritamente as restrições de capacidade ($0 \le f(u,v) \le c(u,v)$), mas que permite que a quantidade total de fluxo que ingressa em um vértice $v \in V \setminus \{s, t\}$ seja maior ou igual à quantidade que dele diverge. A diferença algébrica entre o influxo e o efluxo de um vértice é denominada excesso, denotada por $e(v)$:
\begin{equation}
    e(v) = \sum_{\{u \mid (u,v) \in A\}} f(u,v) - \sum_{\{w \mid (v,w) \in A\}} f(v,w) \ge 0
\end{equation}

Um vértice intermediário que possui excesso estritamente positivo ($e(v) > 0$) é classificado como um vértice ativo. A premissa central dos algoritmos de pré-fluxo consiste em movimentar localmente este excesso pela rede até que ele alcance o sumidouro $t$ ou, em caso de estrangulamento de capacidade, seja devolvido à fonte $s$ (veja~\cite{cook1998combinatorial}).

Para garantir que o excesso não circule infinitamente em ciclos residuais e que convirja para os terminais, introduz-se uma função de rótulos de distância, classicamente interpretada como uma função de altura (\textit{height function}) $h: V \rightarrow \mathbb{N}$. A mecânica do algoritmo é frequentemente comparada a uma malha topográfica de fluidos sujeita à gravidade: o fluxo só pode ser roteado de um vértice $u$ para um vértice $v$ se o arco possuir capacidade residual ($c_f(u,v) > 0$) e se a origem estiver topologicamente posicionada acima do destino, caracterizando um aclive descendente.

Uma função de altura é considerada matematicamente válida para um pré-fluxo se satisfizer três condições peremptórias (veja~\cite{korte2018combinatorial}):
\begin{enumerate}
    \item A altura do sumidouro é nula: $h(t) = 0$.
    \item A altura da fonte é fixada na cardinalidade do digrafo para atestar o isolamento: $h(s) = |V|$.
    \item Para todo arco residual viável $(u,v) \in A_f$, a inclinação é controlada: $h(u) \le h(v) + 1$.
\end{enumerate}

A terceira condição garante que o excesso não despencará em "quedas livres" descontroladas e, principalmente, que não existe nenhum caminho direcionado viável de $s$ para $t$ na rede residual, atestando a otimalidade topológica do sistema a todo instante. O algoritmo termina quando o conjunto de vértices ativos torna-se vazio ($e(v) = 0, \forall v \in V \setminus \{s, t\}$). Neste ponto exato, o pré-fluxo converte-se em um fluxo viável e, pela ausência de caminhos da fonte ao sumidouro, o Teorema do Fluxo Máximo-Corte Mínimo certifica que este fluxo é máximo.

\subsubsection{A Família Push-Relabel}

A mecânica operacional da família de algoritmos \textit{Push-Relabel}, concebida originalmente por Goldberg e Tarjan (veja~\cite{goldberg1988new}), assenta-se sobre a manipulação iterativa de vértices ativos visando o esgotamento topológico de seus excessos. A inicialização do sistema satura artificialmente todos os arcos originais adjacentes à fonte $s$, injetando um pré-fluxo maciço na rede. Simultaneamente, a altura da fonte é fixada em $h(s) = |V|$ e a dos demais vértices é inicializada em $h(v) = 0$, estabelecendo um desnível rigoroso que força o fluido a buscar o sumidouro (veja~\cite{ahuja1993network}).

Durante o ciclo de resolução principal, o algoritmo inspeciona vértices ativos ($e(u) > 0$) e aplica fundamentalmente duas operações elementares para movimentar o pré-fluxo pela rede residual:

\begin{itemize}
    \item \textbf{Operação de Empurrar (\textit{Push}):} Esta ação é executada se, e somente se, o vértice ativo $u$ possuir um arco residual $(u,v)$ estritamente admissível. Um arco é considerado admissível quando apresenta capacidade residual não-nula ($c_f(u,v) > 0$) e um declive topológico unitário perfeito em direção ao destino ($h(u) = h(v) + 1$). A quantidade de fluxo transferida é algebricamente delimitada por $\delta = \min(e(u), c_f(u,v))$. Caso o empurrão transfira exatamente $c_f(u,v)$, o arco é plenamente saturado e removido temporariamente de $A_f$; caso contrário, trata-se de um empurrão não-saturante, cujo efeito primário e imediato é zerar o excesso do vértice $u$ (veja~\cite{cook1998combinatorial}).

    \item \textbf{Operação de Rerrotular (\textit{Relabel}):} Quando um vértice ativo $u$ encontra-se topologicamente "preso" — isto é, possui excesso de recursos, mas nenhuma de suas conexões com capacidade residual aponta para um vizinho em um nível estritamente inferior —, a operação de \textit{Push} torna-se inviável. Para romper este impasse estrutural, o algoritmo atualiza a altura do vértice para $h(u) = 1 + \min \{h(v) \mid c_f(u,v) > 0\}$. Esta elevação pontual garante a recriação de um declive geométrico válido, permitindo que o excesso retome sua movimentação iterativa (veja~\cite{korte2018combinatorial}).
\end{itemize}

A eficiência global do algoritmo é intrinsecamente dependente da política de seleção: qual vértice ativo deve ser processado primeiro? A implementação canônica emprega uma estrutura de dados de fila (\textit{Queue}) para gerenciar estes vértices sob uma política \textit{First-In, First-Out} (FIFO). Demonstra-se matematicamente que a ordem de processamento FIFO amortece severamente o número de empurrões não-saturantes necessários para a convergência. Graças a esta amortização estrutural, o limite assintótico de tempo de pior caso para o algoritmo Push-Relabel padrão é cravado analiticamente em $\mathcal{O}(|V|^3)$ (veja~\cite{goldberg1988new}). A materialização em C++ deste algoritmo, herdando a arquitetura segura da interface polimórfica, encontra-se detalhada no Código~\ref{lst:push_relabel} (disponível no Apêndice~\ref{ap:push_relabel}).

\subsubsection{Otimização Empírica: A Heurística de Lacuna (\textit{Gap Heuristic})}

Para suprimir este gargalo crítico, a literatura moderna incorpora uma modificação estrutural de estado profundo conhecida como Heurística de Lacuna (\textit{Gap Heuristic}). A arquitetura do sistema é expandida para manter um histograma dinâmico (\textit{array} de frequências) que registra em tempo constante o número exato de vértices presentes em cada estrato de altura da rede.

A fundamentação desta heurística baseia-se na observação de descontinuidades no histograma de alturas. Se, após uma operação de \textit{relabel}, um determinado nível topológico $k$ (onde $0 < k < |V|$) esvaziar-se completamente, forma-se uma lacuna (\textit{gap}). A existência desta lacuna implica que o sumidouro $t$ tornou-se inalcançável para todo vértice com altura estritamente superior a $k$.

A heurística intervém na rede ao detectar a lacuna, elevando a altura de todos os vértices acima de $k$ para o patamar de isolamento ($|V| + 1$) (veja~\cite{korte2018combinatorial}). Esta operação reduz o número de operações de \textit{relabel} redundantes, otimizando a convergência do algoritmo em grafos densos (veja~\cite{goldberg1988new}).

\subsection{Reduções e Aplicações Combinatórias}
O Problema do Fluxo Máximo é frequentemente utilizado como base para resolver outros problemas de otimização combinatória via redução polinomial (veja~\cite{cook1998combinatorial}). O processo consiste em mapear a instância de um problema original para uma rede de fluxo equivalente, computar o fluxo máximo e traduzir o estado da rede residual de volta para a solução do problema original.

Um digrafo $D = (V, A)$ é classificado como bipartido se o seu conjunto de vértices $V$ puder ser particionado em dois subconjuntos disjuntos independentes, convencionalmente denotados por $L$ (esquerdo) e $R$ (direito), de modo que todo arco $a \in A$ conecte obrigatoriamente um vértice de $L$ a um vértice de $R$. Esta topologia é a estrutura canônica para a modelagem de sistemas de alocação bidirecional.

\subsubsection{Emparelhamento Bipartido Máximo}

O Problema do Emparelhamento Bipartido de Cardinalidade Máxima (\textit{Maximum Bipartite Matching}) busca encontrar o maior subconjunto de arcos $M \subseteq A$ tal que nenhum vértice do digrafo incida em mais de um arco de $M$. Em termos práticos, representa a maximização de pares exclusivos entre os conjuntos $L$ e $R$.

Embora existam algoritmos específicos para este fim, o problema pode ser elegantemente reduzido ao cômputo de um fluxo máximo (veja~\cite{ahuja1993network}). A redução algorítmica procede através da construção de uma rede de fluxo direcionada $D' = (V', A')$ baseada no digrafo original $D$:
\begin{enumerate}
    \item Adicionam-se dois vértices artificiais ao sistema: uma super-fonte $s$ e um super-sumidouro $t$.
    \item Para cada vértice $u \in L$, cria-se um arco direcionado $(s, u)$ com capacidade exatamente igual a 1.
    \item Para cada vértice $v \in R$, cria-se um arco direcionado $(v, t)$ com capacidade exatamente igual a 1.
    \item Todos os arcos originais do digrafo não-direcionado são orientados de $L$ para $R$, recebendo capacidade infinita (ou alternativamente, capacidade 1).
\end{enumerate}

A prova de correção desta modelagem repousa sobre o Teorema da Integralidade (veja~\cite{korte2018combinatorial}). Como todas as capacidades do subdigrafo injetor ($s \rightarrow L$) e extrator ($R \rightarrow t$) são inteiras (estritamente 1), os algoritmos de fluxo garantem matematicamente o retorno de um vetor de fluxo viável inteiramente composto por zeros e uns.

Consequentemente, a restrição de capacidade unitária nos arcos conectados à super-fonte e ao super-sumidouro força a lei de conservação a atuar como um restritor de exclusividade: se $1$ unidade de fluxo transita por um vértice $u \in L$, ela só pode ser roteada por um único arco em direção a $R$. O valor numérico do fluxo máximo computado pela rede $D'$ equivale perfeitamente à cardinalidade $|M|$ do emparelhamento máximo, e os arcos ativos ($f(u,v) = 1$) que cruzam de $L$ para $R$ constituem os pares da solução exata.

\subsubsection{O Teorema de König e a Cobertura de Vértices Mínima}

O Problema da Cobertura de Vértices Mínima (\textit{Minimum Vertex Cover} - MVC) consiste em determinar o menor subconjunto de vértices $C \subseteq V$ tal que todo arco do digrafo incida em pelo menos um vértice pertencente a $C$. Em digrafos arbitrários, este problema é classificado como NP-Difícil. Contudo, quando a topologia do digrafo é estritamente bipartida, o problema pode ser resolvido em tempo polinomial através da sua profunda relação dual com o emparelhamento máximo.

Esta relação de dualidade é formalizada pelo Teorema de König (veja~\cite{korte2018combinatorial}), que estabelece que, em qualquer digrafo bipartido, a cardinalidade do emparelhamento máximo é matematicamente idêntica à cardinalidade da cobertura de vértices mínima ($|M| = |C|$). Sob a ótica da teoria de redes, o Teorema de König nada mais é do que uma manifestação estrutural do Teorema do Fluxo Máximo-Corte Mínimo de Ford-Fulkerson aplicado à modelagem com super-fonte e super-sumidouro (veja~\cite{ahuja1993network}).

A extração computacional dos vértices que compõem a cobertura mínima exige uma análise minuciosa da rede residual $A_f$ gerada logo após a convergência do fluxo máximo. A biblioteca desenvolvida neste projeto recupera o conjunto $C$ através do seguinte procedimento algorítmico baseado em cortes mínimos:

\begin{enumerate}
    \item Após computar o emparelhamento máximo, executa-se uma Busca em Largura (BFS) ou Busca em Profundidade (DFS) partindo da super-fonte $s$ na rede residual $A_f$, transitando exclusivamente por arcos que possuam capacidade residual estritamente positiva ($c_f(u,v) > 0$).
    \item Esta travessia particiona os vértices do digrafo bipartido original em dois conjuntos disjuntos: o conjunto dos vértices alcançáveis ($S$) e o conjunto dos vértices inalcançáveis ($T$). Este particionamento define o corte mínimo da rede.
    \item A Cobertura de Vértices Mínima $C$ é então construída selecionando-se os vértices da partição esquerda original que não foram alcançados pela busca ($L \setminus S$) unidos aos vértices da partição direita que foram alcançados ($R \cap S$). Formalmente: $C = (L \cap T) \cup (R \cap S)$.
\end{enumerate}

A solidez desta extração reside no fato de que, pelas propriedades de conservação do corte mínimo em capacidades unitárias, não pode existir nenhum arco residual viável cruzando de $S$ para $T$. Consequentemente, a seleção $(L \cap T) \cup (R \cap S)$ garante de forma inequívoca que todos os arcos originais do digrafo bipartido estão "cobertos" por, no mínimo, um destes vértices selecionados, consolidando a prova empírica do Teorema de König.

\subsubsection{Conjunto Independente Máximo}

O Problema do Conjunto Independente Máximo (\textit{Maximum Independent Set} - MIS) visa identificar o maior subconjunto de vértices $I \subseteq V$ no qual nenhuma dupla de vértices seja adjacente entre si. Em outras palavras, nenhum arco do digrafo pode possuir ambas as extremidades contidas no conjunto $I$.

Assim como a Cobertura de Vértices, este problema é NP-Difícil para digrafos gerais, mas admite uma solução polinomial elegante e direta em digrafos bipartidos graças à dualidade complementar. A relação matemática entre a Cobertura de Vértices Mínima ($C$) e o Conjunto Independente Máximo ($I$) é inversamente proporcional e estritamente complementar: se um conjunto de vértices $C$ cobre todos os arcos de um digrafo, a remoção de $C$ elimina todas as conexões possíveis entre os vértices restantes (veja~\cite{cook1998combinatorial}).

Portanto, os vértices que sobrevivem a esta remoção não possuem arcos entre si, formando, por definição, um conjunto independente. Como $C$ foi minimizado pelo Teorema de König e pelo algoritmo de Corte Mínimo, o conjunto complementar resultante é matematicamente garantido como sendo o máximo possível.

Desta forma, a extração computacional do Conjunto Independente Máximo dispensa a necessidade de novos roteamentos de fluxo ou buscas em digrafos residuais. A biblioteca desenvolvida computa o conjunto $I$ através de uma operação trivial de diferença de conjuntos em tempo $\mathcal{O}(|V|)$:
\begin{equation}
    I = V \setminus C
\end{equation}

A cardinalidade do Conjunto Independente Máximo é dada algebricamente por $|I| = |V| - |C|$. Esta sequência de reduções demonstra a aplicabilidade do modelo de fluxo em redes: a partir de uma única construção baseada em super-fonte e super-sumidouro, o sistema soluciona simultaneamente o Emparelhamento Máximo, a Cobertura de Vértices Mínima e o Conjunto Independente Máximo.

\subsection{Conectividade e Roteamento}

A análise de conectividade de um digrafo é um pilar fundamental no design de redes de telecomunicações, logística e tolerância a falhas. A premissa central é determinar a resiliência estrutural de um sistema: quantas rotas independentes existem para transmitir informações ou recursos de um ponto de origem $s$ a um destino $t$ de forma que, se um componente falhar, as demais rotas permaneçam intactas? O cômputo do fluxo máximo fornece o ferramental matemático exato para responder a esta indagação sob duas perspectivas restritivas.

\subsubsection{Caminhos Disjuntos por Arcos e por Vértices}

A formulação mais elementar da independência de rotas é o Problema dos Caminhos Disjuntos por Arcos (\textit{Edge-Disjoint Paths}). Dois caminhos são considerados disjuntos por arcos se não compartilharem nenhuma conexão física em suas trajetórias, embora possam cruzar-se nos mesmos vértices intermediários.

A redução deste problema computacional para o modelo de fluxo em redes é direta. Dado um digrafo direcionado $D = (V, A)$, constrói-se uma rede de fluxo onde a restrição de capacidade de absolutamente todos os arcos é fixada em um valor unitário ($c(u,v) = 1, \forall (u,v) \in A$). Pela aplicação direta do Teorema da Integralidade (veja~\cite{cook1998combinatorial}), o algoritmo de fluxo máximo garantirá que cada arco transporte, no máximo, $1$ unidade de fluxo. Consequentemente, o valor absoluto do fluxo computado será numericamente equivalente ao número exato de caminhos disjuntos por arcos existentes entre $s$ e $t$.

Uma restrição topológica consideravelmente mais severa é imposta pelo Problema dos Caminhos Disjuntos por Vértices (\textit{Vertex-Disjoint Paths}). Neste cenário, os caminhos não podem compartilhar sequer um vértice intermediário (cruzamento), garantindo isolamento total. Como os algoritmos canônicos de fluxo operam controlando as capacidades dos arcos e não dos vértices, a modelagem exige uma transformação geométrica avançada conhecida como Divisão de Vértices (\textit{Vertex Splitting}) (veja~\cite{ahuja1993network}).

A técnica de \textit{Vertex Splitting} altera a topologia do digrafo original $D$ para gerar uma nova rede de fluxo estendida $A'$ através da seguinte heurística de mapeamento:
\begin{enumerate}
    \item Cada vértice original $v \in V \setminus \{s, t\}$ é bipartido em dois vértices artificiais distintos: um vértice de entrada ($v_{in}$) e um vértice de saída ($v_{out}$).
    \item Um arco direcionado interno é criado conectando as duas metades, $(v_{in}, v_{out})$, e recebe uma capacidade estritamente unitária: $c(v_{in}, v_{out}) = 1$.
    \item Todos os arcos que originalmente apontavam para $v$ no digrafo $D$ são redirecionados para apontar exclusivamente para $v_{in}$.
    \item Todos os arcos que originalmente divergiam de $v$ no digrafo $D$ passam a emergir exclusivamente a partir de $v_{out}$ (com capacidade 1 ou infinita, dependendo da formulação).
\end{enumerate}

Esta transformação impõe uma restrição de capacidade nos vértices: ao forçar todo o fluxo que transita pelo nó original $v$ a atravessar um arco interno $(v_{in}, v_{out})$ de capacidade $1$, a rede garante que não mais do que $1$ unidade de fluxo possa cruzar aquele vértice. Ao submeter esta nova topologia estendida $D'$ aos algoritmos implementados, o fluxo máximo computado determina o número máximo de rotas disjuntas por vértices (veja~\cite{korte2018combinatorial}).

\subsubsection{O Teorema de Menger e a Dualidade da Conectividade}

A determinação empírica do número de caminhos disjuntos resolve o aspecto do roteamento, mas a análise de vulnerabilidade da rede exige a identificação precisa dos pontos críticos de falha. A teoria que unifica a maximização de rotas independentes com a minimização de desconexões foi formulada por Karl Menger em 1927, estabelecendo um dos pilares mais elegantes da Otimização Combinatória.

O Teorema de Menger apresenta duas formulações complementares que espelham perfeitamente os cenários de arcos e vértices (veja~\cite{cook1998combinatorial}). Na sua versão por arcos, o teorema postula que, em qualquer digrafo direcionado, o número máximo de caminhos disjuntos por arcos entre dois vértices $s$ e $t$ é exatamente igual ao número mínimo de arcos que precisam ser removidas para desconectar todos os caminhos direcionados de $s$ para $t$ (o Corte Mínimo de Arcos).

De forma análoga, na sua versão por vértices, o teorema afirma que o número máximo de caminhos internamente disjuntos por vértices entre $s$ e $t$ é rigorosamente idêntico ao número mínimo de vértices intermediários cuja remoção fragmenta a rede e isola o sumidouro da fonte (o Corte Mínimo de Vértices).

Sob a perspectiva da engenharia de algoritmos desenvolvida neste projeto, o Teorema de Menger é, fundamentalmente, um corolário estrutural do Teorema do Fluxo Máximo-Corte Mínimo de Ford-Fulkerson aplicado a redes de capacidade unitária (veja~\cite{ahuja1993network}). A prova de que estas dualidades são faces da mesma moeda matemática é confirmada pelo método de extração do corte: para identificar quais arcos ou vértices específicos devem ser removidos para causar a falha da rede, a biblioteca reutiliza a mesma rotina de Busca em Largura (BFS) na rede residual $A_f$ empregada no Teorema de König.

No caso do \textit{Vertex Splitting}, a BFS varre o digrafo residual particionando os vértices expandidos. Os arcos artificiais internos $(v_{in}, v_{out})$ que conectam o conjunto alcançável $S$ ao conjunto inalcançável $T$ no corte mínimo apontam inequivocamente para os vértices originais $v \in V$ que compõem o gargalo estrutural da topologia (veja~\cite{korte2018combinatorial}).

Estes resultados confirmam a versatilidade do modelo de fluxo em redes. Através de abstrações adequadas e da implementação orientada a objetos, o núcleo algorítmico demonstra a capacidade de resolver problemas que variam desde a alocação de recursos em grafos bipartidos até a análise de conectividade topológica.

\section{Conclusões preliminares e próximos passos}

Os resultados parciais obtidos até o presente momento em cumprimento às Atividades 1 a 5 do cronograma demonstram o sucesso na consolidação de uma infraestrutura computacional robusta e teoricamente fundamentada para a resolução do Problema do Fluxo Máximo. A adoção de uma arquitetura Orientada a Objetos em C++, ancorada pelos padrões \textit{Factory Method} e \textit{Prototype}, provou-se altamente eficaz para o encapsulamento da topologia das redes residuais.

Sob a ótica algorítmica, o projeto materializou com êxito os métodos clássicos propostos na literatura. A implementação avançou desde o método de Edmonds-Karp, passando pela otimização em fases com o Digrafo de Níveis de Dinic, até atingir o estado da arte com a família de prefluxos através do algoritmo \textit{Push-Relabel} (Goldberg-Tarjan). A correção matemática destas implementações foi validada através da sua aplicação bem-sucedida em técnicas de modelagem geométrica, solucionando problemas complexos como o Emparelhamento Bipartido Máximo e o empacotamento de Caminhos Disjuntos (por arcos e vértices).

Tendo consolidado a teoria e a base de código para o fluxo máximo, e em conformidade com o cronograma estabelecido (Atividade 6), as próximas etapas da Iniciação Científica (Atividades 7 a 11) concentrar-se-ão na expansão do escopo teórico e na validação empírica rigorosa, contemplando os seguintes passos:

\begin{itemize}
    \item \textbf{Estudo e Implementação do Fluxo de Custo Mínimo (Atividades 7 e 8):} A pesquisa transicionará para o problema de fluxo de custo mínimo, que generaliza as abordagens atuais ao incorporar restrições financeiras/logísticas ao transporte. O foco teórico e prático será o estudo e a implementação em C++ do algoritmo clássico \textit{network simplex}.
    \item \textbf{Testes Experimentais (Atividade 9):} Condução de baterias de testes utilizando as implementações desenvolvidas. Para garantir o rigor científico e a reprodutibilidade, os experimentos serão executados sobre bases de dados e instâncias de \textit{benchmarks} amplamente reconhecidas na literatura, em especial as coleções padronizadas do DIMACS.
    \item \textbf{Análise Comparativa de Desempenho (Atividade 10):} Execução de um \textit{benchmarking} intensivo para avaliar o tempo de execução e o consumo de memória dos algoritmos implementados. Esta análise contrastará o comportamento empírico dos métodos de fluxo máximo (Dinic vs. Push-Relabel) e avaliará a eficiência do \textit{network simplex} sob diferentes topologias de digrafos.
    \item \textbf{Elaboração do Relatório Final (Atividade 11):} Compilação de todos os dados experimentais, análises assintóticas e códigos desenvolvidos em um documento final que sintetize as contribuições do discente para a área de Otimização Combinatória.
\end{itemize}

Desta forma, a infraestrutura estabelecida nesta primeira fase oferece total segurança técnica para as implementações avançadas e baterias de testes que ditarão o ritmo e a conclusão da pesquisa.

\newpage

%\bibliographystyle{plain}{abntex2-al}
\bibliographystyle{plain}
\bibliography{cit}

\newpage
\appendix

\begin{center}{\Large\textbf{Apêndice}}\end{center}
\section{Implementação da Interface Base ``Flow-Network''}
\label{ap:flownetwork}

\begin{lstlisting}[caption={Implementação em C++ da interface base \lstinline{FlowNetwork} para a modelagem abstrata de redes de fluxo.}, label={lst:flownetwork}]
#include <algorithm>
#include <iostream>
#include <limits>
#include <memory>
#include <queue>
#include <vector>

using Long = long long;
using Size = std::size_t;

constexpr Long INF = std::numeric_limits<Long>::max() >> 8;
constexpr Size MAX = std::numeric_limits<Size>::max() >> 8;

class FlowNetwork
{
public:
	struct Edge
	{
		Size from, to;
		Long capacity, flow;
	};

	explicit FlowNetwork(const Size n) : size(n), adj(n) {}
	virtual ~FlowNetwork() = default;

	virtual std::unique_ptr<FlowNetwork> make(const Size n) const = 0;
	virtual std::unique_ptr<FlowNetwork> clone() const = 0;

	virtual void add_edge(
	    const Size from, const Size to, const Long capacity,
	    const Long reverse_capacity = 0
	)
	{
		adj[from].push_back(edges.size());
		edges.push_back({from, to, capacity, 0});
		adj[to].push_back(edges.size());
		edges.push_back({to, from, reverse_capacity, 0});
	}

	virtual Long compute_max_flow(const Size source, const Size sink) = 0;

	[[nodiscard]] Size get_size() const
	{
		return size;
	}

	[[nodiscard]] const std::vector<Edge> &get_edges() const
	{
		return edges;
	}

	[[nodiscard]] const std::vector<std::vector<Size>> &get_adj() const
	{
		return adj;
	}

protected:
	Size size;
	std::vector<Edge> edges;
	std::vector<std::vector<Size>> adj;

	[[nodiscard]] Long get_residual_capacity(const Size edge_id) const
	{
		return edges[edge_id].capacity - edges[edge_id].flow;
	}

	void push_flow(const Size edge_id, const Long flow_amount)
	{
		edges[edge_id].flow += flow_amount;
		edges[edge_id ^ 1ULL].flow -= flow_amount;
	}
};
\end{lstlisting}

\newpage
\section{Implementação da Classe ``Edmonds-Karp''}
\label{ap:edmonds_karp}

\begin{lstlisting}[caption={Implementação em C++ da classe \lstinline{EdmondsKarp}, baseada na seleção iterativa de caminhos mais curtos via Busca em Largura (BFS).}, label={lst:edmonds_karp}]
class EdmondsKarp : public FlowNetwork
{
public:
	explicit EdmondsKarp(const Size n) : FlowNetwork(n) {}

	static std::unique_ptr<FlowNetwork> create(const Size n)
	{
		return std::make_unique<EdmondsKarp>(n);
	}

	std::unique_ptr<FlowNetwork> make(const Size n) const override
	{
		return std::make_unique<EdmondsKarp>(n);
	}

	std::unique_ptr<FlowNetwork> clone() const override
	{
		return std::make_unique<EdmondsKarp>(*this);
	}

	Long compute_max_flow(const Size source, const Size sink) override
	{
		Long total_flow = 0;
		Long new_flow = 0;
		std::vector<Size> parent_edge(size);

		while ((new_flow = bfs(source, sink, parent_edge)) > 0)
		{
			total_flow += new_flow;
			update_path_flow(source, sink, parent_edge, new_flow);
		}

		return total_flow;
	}

private:
	Long bfs(const Size source, const Size sink, std::vector<Size> &parent_edge)
	{
		std::fill(parent_edge.begin(), parent_edge.end(), MAX);
		std::queue<std::pair<Size, Long>> queue;

		parent_edge[source] = source;
		queue.push({source, INF});

		while (!queue.empty())
		{
			const auto [current_node, current_flow] = queue.front();
			queue.pop();

			for (const Size edge_id : adj[current_node])
			{
				const Size next_node = edges[edge_id].to;
				const Long residual_capacity = get_residual_capacity(edge_id);

				if (parent_edge[next_node] != MAX || residual_capacity <= 0)
					continue;

				parent_edge[next_node] = edge_id;
				const Long pushed_flow = std::min(current_flow, residual_capacity);

				if (next_node == sink)
					return pushed_flow;

				queue.push({next_node, pushed_flow});
			}
		}
		return 0;
	}

	void update_path_flow(
	    const Size source, const Size sink, const std::vector<Size> &parent_edge,
	    const Long new_flow
	)
	{
		Size current_node = sink;
		while (current_node != source)
		{
			const Size edge_id = parent_edge[current_node];
			push_flow(edge_id, new_flow);
			current_node = edges[edge_id].from;
		}
	}
};
\end{lstlisting}

\newpage
\section{Implementação da Classe ``Dinic''}
\label{ap:dinic}

\begin{lstlisting}[caption={Implementação em C++ da classe \lstinline{Dinic}, estruturada com Digrafo de Níveis e otimização de ponteiros mortos no cálculo do fluxo bloqueador.}, label={lst:dinic}]
class Dinic : public FlowNetwork
{
public:
	explicit Dinic(const Size n) : FlowNetwork(n), level(n), next_edge_ptr(n) {}

	static std::unique_ptr<FlowNetwork> create(const Size n)
	{
		return std::make_unique<Dinic>(n);
	}

	std::unique_ptr<FlowNetwork> make(const Size n) const override
	{
		return std::make_unique<Dinic>(n);
	}

	std::unique_ptr<FlowNetwork> clone() const override
	{
		return std::make_unique<Dinic>(*this);
	}

	Long compute_max_flow(const Size source, const Size sink) override
	{
		Long total_flow = 0;
		while (bfs(source, sink))
		{
			std::fill(next_edge_ptr.begin(), next_edge_ptr.end(), 0);
			while (const Long flow_pushed = dfs(source, sink, INF))
			{
				total_flow += flow_pushed;
			}
		}
		return total_flow;
	}

private:
	std::vector<Size> level;
	std::vector<Size> next_edge_ptr;

	bool bfs(const Size source, const Size sink)
	{
		std::fill(level.begin(), level.end(), MAX);
		std::queue<Size> queue;

		level[source] = 0;
		queue.push(source);

		while (!queue.empty())
		{
			const Size current_node = queue.front();
			queue.pop();

			for (const Size edge_id : adj[current_node])
			{
				const Size next_node = edges[edge_id].to;
				const Long residual_capacity = get_residual_capacity(edge_id);

				if (level[next_node] != MAX || residual_capacity <= 0)
					continue;

				level[next_node] = level[current_node] + 1;
				queue.push(next_node);
			}
		}
		return level[sink] != MAX;
	}

	Long dfs(const Size current_node, const Size sink, const Long flow_pushed)
	{
		if (flow_pushed == 0 || current_node == sink)
			return flow_pushed;

		for (Size &ptr = next_edge_ptr[current_node]; ptr < adj[current_node].size();
		     ++ptr)
		{
			const Size edge_id = adj[current_node][ptr];
			const Size next_node = edges[edge_id].to;
			const Long residual_capacity = get_residual_capacity(edge_id);

			if (level[current_node] + 1 != level[next_node] ||
			    residual_capacity <= 0)
				continue;

			const Long bottleneck = std::min(flow_pushed, residual_capacity);
			const Long flow_transmitted = dfs(next_node, sink, bottleneck);

			if (flow_transmitted == 0)
				continue;

			push_flow(edge_id, flow_transmitted);
			return flow_transmitted;
		}
		return 0;
	}
};
\end{lstlisting}

\newpage
\section{Implementação da Classe ``Push-Relabel''}
\label{ap:push_relabel}

\begin{lstlisting}[caption={Implementação em C++ da classe \lstinline{PushRelabel}, baseada em pré-fluxos topológicos e política de processamento em fila (FIFO).}, label={lst:push_relabel}]
class PushRelabel : public FlowNetwork
{
public:
	explicit PushRelabel(const Size n)
	    : FlowNetwork(n), height(n), excess(n), next_edge_ptr(n)
	{
	}

	static std::unique_ptr<FlowNetwork> create(const Size n)
	{
		return std::make_unique<PushRelabel>(n);
	}

	std::unique_ptr<FlowNetwork> make(const Size n) const override
	{
		return std::make_unique<PushRelabel>(n);
	}

	std::unique_ptr<FlowNetwork> clone() const override
	{
		return std::make_unique<PushRelabel>(*this);
	}

	Long compute_max_flow(const Size source, const Size sink) override
	{
		std::queue<Size> active_queue;
		std::vector<bool> in_queue(size, false);

		initialize_preflow(source, sink, active_queue, in_queue);

		while (!active_queue.empty())
		{
			const Size current_node = active_queue.front();
			active_queue.pop();
			in_queue[current_node] = false;
			discharge_node(current_node, active_queue, in_queue, source, sink);
		}

		return excess[sink];
	}

private:
	std::vector<Size> height;
	std::vector<Long> excess;
	std::vector<Size> next_edge_ptr;

	void initialize_preflow(
	    const Size source, const Size sink, std::queue<Size> &active_queue,
	    std::vector<bool> &in_queue
	)
	{
		std::fill(height.begin(), height.end(), 0);
		std::fill(excess.begin(), excess.end(), 0);
		std::fill(next_edge_ptr.begin(), next_edge_ptr.end(), 0);

		height[source] = size;
		excess[source] = INF;

		for (const Size edge_id : adj[source])
		{
			if (get_residual_capacity(edge_id) <= 0)
			{
				continue;
			}

			push_preflow(source, edge_id);
			const Size next_node = edges[edge_id].to;

			if (next_node == source || next_node == sink || in_queue[next_node])
			{
				continue;
			}

			active_queue.push(next_node);
			in_queue[next_node] = true;
		}
	}

	void discharge_node(
	    const Size current_node, std::queue<Size> &active_queue,
	    std::vector<bool> &in_queue, const Size source, const Size sink
	)
	{
		while (excess[current_node] > 0)
		{
			if (next_edge_ptr[current_node] >= adj[current_node].size())
			{
				relabel_node(current_node);
				next_edge_ptr[current_node] = 0;
				continue;
			}

			const Size edge_id = adj[current_node][next_edge_ptr[current_node]];
			const Size next_node = edges[edge_id].to;
			const Long residual_capacity = get_residual_capacity(edge_id);

			if (residual_capacity > 0 &&
			    height[current_node] == height[next_node] + 1)
			{
				push_preflow(current_node, edge_id);

				if (next_node != source && next_node != sink &&
				    !in_queue[next_node] && excess[next_node] > 0)
				{
					active_queue.push(next_node);
					in_queue[next_node] = true;
				}
			}
			else
			{
				next_edge_ptr[current_node]++;
			}
		}
	}

	void relabel_node(const Size current_node)
	{
		Size min_height = MAX;
		for (const Size edge_id : adj[current_node])
		{
			if (get_residual_capacity(edge_id) <= 0)
			{
				continue;
			}
			min_height = std::min(min_height, height[edges[edge_id].to]);
		}

		if (min_height != MAX)
		{
			height[current_node] = min_height + 1;
		}
	}

	void push_preflow(const Size current_node, const Size edge_id)
	{
		const Size next_node = edges[edge_id].to;
		const Long residual_capacity = get_residual_capacity(edge_id);
		const Long flow_to_push = std::min(excess[current_node], residual_capacity);

		if (flow_to_push == 0)
			return;

		push_flow(edge_id, flow_to_push);
		excess[current_node] -= flow_to_push;
		excess[next_node] += flow_to_push;
	}
};
\end{lstlisting}

\newpage
\section{Implementação da Classe ``Push-Relabel Improved''}
\label{ap:push_relabel_improved}

\begin{lstlisting}[caption={Implementação em C++ da classe \lstinline{PushRelabelImproved}, otimizada empiricamente com a Heurística de Lacuna (\textit{Gap Heuristic}).}, label={lst:push_relabel_improved}]
class PushRelabelImproved : public FlowNetwork
{
public:
	explicit PushRelabelImproved(const Size n) : FlowNetwork(n) {}

	static std::unique_ptr<FlowNetwork> create(const Size n)
	{
		return std::make_unique<PushRelabelImproved>(n);
	}

	std::unique_ptr<FlowNetwork> make(const Size n) const override
	{
		return std::make_unique<PushRelabelImproved>(n);
	}

	std::unique_ptr<FlowNetwork> clone() const override
	{
		return std::make_unique<PushRelabelImproved>(*this);
	}

	Long compute_max_flow(const Size source, const Size sink) override
	{
		initialize_preflow(source, sink);

		while (true)
		{
			if (active_buckets[highest_active].empty())
			{
				if (highest_active == 0)
					break;
				highest_active--;
				continue;
			}

			const Size current_node = active_buckets[highest_active].back();
			active_buckets[highest_active].pop_back();

			if (height[current_node] != highest_active)
			{
				continue;
			}

			discharge_node(current_node, source, sink);
		}

		return excess[sink];
	}

private:
	std::vector<Size> height;
	std::vector<Long> excess;
	std::vector<Size> height_count;
	std::vector<Size> next_edge_ptr;
	std::vector<std::vector<Size>> active_buckets;
	Size highest_active;

	void initialize_preflow(const Size source, const Size sink)
	{
		const Size max_height_limit = size * 2 + 1;

		height.assign(size, 0);
		excess.assign(size, 0);
		height_count.assign(max_height_limit, 0);
		next_edge_ptr.assign(size, 0);
		active_buckets.assign(max_height_limit, std::vector<Size>());
		highest_active = 0;

		height[source] = size;
		height_count[size] = 1;
		height_count[0] = size - 1;

		excess[source] = INF;
		for (const Size edge_id : adj[source])
		{
			push_preflow(source, edge_id, source, sink);
		}
	}

	void discharge_node(const Size current_node, const Size source, const Size sink)
	{
		while (excess[current_node] > 0)
		{
			if (next_edge_ptr[current_node] >= adj[current_node].size())
			{
				relabel_node(current_node, source, sink);
				next_edge_ptr[current_node] = 0;
				continue;
			}

			const Size edge_id = adj[current_node][next_edge_ptr[current_node]];
			const Size next_node = edges[edge_id].to;
			const Long residual_capacity = get_residual_capacity(edge_id);

			if (residual_capacity > 0 &&
			    height[current_node] == height[next_node] + 1)
			{
				push_preflow(current_node, edge_id, source, sink);
			}
			else
			{
				next_edge_ptr[current_node]++;
			}
		}
	}

	void relabel_node(const Size current_node, const Size source, const Size sink)
	{
		Size min_height = MAX;
		for (const Size edge_id : adj[current_node])
		{
			if (get_residual_capacity(edge_id) <= 0)
			{
				continue;
			}

			min_height = std::min(min_height, height[edges[edge_id].to]);
		}

		if (min_height == MAX)
		{
			return;
		}

		const Size old_height = height[current_node];
		height_count[old_height]--;

		height[current_node] = min_height + 1;
		height_count[height[current_node]]++;
		enqueue_active_node(current_node, source, sink);

		if (height_count[old_height] == 0 && old_height < size)
		{
			apply_gap_heuristic(old_height, source, sink);
		}
	}

	void apply_gap_heuristic(
	    const Size gap_height, const Size source, const Size sink
	)
	{
		for (Size i = 0; i < size; ++i)
		{
			if (i == source || height[i] <= gap_height || height[i] >= size)
			{
				continue;
			}

			height_count[height[i]]--;
			height[i] = std::max(height[i], size + 1);
			height_count[height[i]]++;
			enqueue_active_node(i, source, sink);
		}
	}

	void push_preflow(
	    const Size current_node, const Size edge_id, const Size source,
	    const Size sink
	)
	{
		const Size next_node = edges[edge_id].to;
		const Long residual_capacity = get_residual_capacity(edge_id);

		if (residual_capacity == 0 || height[current_node] <= height[next_node])
		{
			return;
		}

		const Long flow_to_push = std::min(excess[current_node], residual_capacity);
		push_flow(edge_id, flow_to_push);
		excess[current_node] -= flow_to_push;
		excess[next_node] += flow_to_push;

		if (excess[next_node] == flow_to_push)
		{
			enqueue_active_node(next_node, source, sink);
		}
	}

	void enqueue_active_node(const Size node, const Size source, const Size sink)
	{
		if (node == source || node == sink || excess[node] <= 0)
		{
			return;
		}

		active_buckets[height[node]].push_back(node);
		highest_active = std::max(highest_active, height[node]);
	}
};
\end{lstlisting}

\newpage
\section{Acesso ao Código-Fonte e Repositório GitHub}
\label{ap:github}

Com o objetivo de garantir a transparência, a reprodutibilidade dos experimentos e facilitar a consulta das implementações discutidas neste relatório, todo o material de desenvolvimento foi organizado em um repositório público na plataforma GitHub.

O repositório \textbf{IC\_Networks\_Flow} contém:
\begin{itemize}
    \item A biblioteca base em C++ com a interface polimórfica \lstinline{FlowNetwork};
    \item Implementações completas dos algoritmos de Edmonds-Karp, Dinic e Push-Relabel (incluindo a \textit{Gap Heuristic});
    \item Scripts para automação de testes e leitura de instâncias no padrão DIMACS;
    \item Documentação técnica sobre a estrutura das classes e instruções para compilação.
\end{itemize}

O código-fonte completo pode ser consultado, clonado e executado através do endereço:

\vspace{0.4cm}
\begin{center}
    \fbox{\url{https://github.com/GabrielFrigo4/IC_Networks_Flow}}
\end{center}
\vspace{0.4cm}

O versionamento do código segue as boas práticas de desenvolvimento de \textit{software}, permitindo o acompanhamento histórico das otimizações realizadas durante cada etapa deste projeto de Iniciação Científica.

\end{document}
